\documentclass[11pt]{article}
\usepackage[margin=2.5cm]{geometry}
\usepackage{amsmath,amssymb,booktabs,microtype}
\usepackage[hidelinks]{hyperref}
\newcommand{\E}{\mathbb{E}}
\newcommand{\Var}{\operatorname{Var}}
\newcommand{\Z}{\mathbb{Z}}
\renewcommand{\thesection}{S\arabic{section}}
\renewcommand{\thetable}{S\arabic{table}}
\renewcommand{\theequation}{S\arabic{equation}}
\title{Supplementary material for\\``Carry Conservation in Approximate Accumulators:\\Quantizer Equivalence and an Entropy Bound''}
\author{Rei Li, Xinyi Li, and Ivy Li}
\date{}

\begin{document}
\maketitle

\noindent Theorem, lemma, equation, section, and table numbers without an
S prefix refer to the main paper. All reported values come from the
study scripts in \texttt{python/accrate}.

\section{Certification of the residue width}

We determine the residue width $k$ of each 8-bit cell by exhausting all
$2^{16}$ operand pairs. For 12- and 16-bit cells, a SAT miter instantiates
two copies of the unmodified RTL. It constrains their low $k$ operand bits
to agree and leaves the high bits free. It then asks whether the signed
errors can differ. An unsatisfiable result proves the residue property
for every high-bit assignment. A concrete pair at $k-1$ establishes
that the width is minimal.

The median SAT proof time was 33--48~ms. Any $w$-bit cell used as the
low segment of a wider accumulator has the property at $k=w$.
The certificate identifies the smallest such $k$. Among 88 nonexact
EvoApproxLib cells, 34 have no residue width $k\le8$ and are excluded.
All are 12- or 16-bit cells, as confirmed by the multiwidth residue-width
certification procedure.

\section{Delayed rounding in AMA5}

With the accumulator state connected as its left operand, AMA5 copies
the fresh residue into the next state, $X_i=R_i$. It emits the
previously stored top bit as carry, $C_i=(X_{i-1})_{k-1}$.
For comparison, the
round-half-up twin has error $E'_n=\sum_i(q(R_i)_{k-1}-R_i)$, so Lemma~2 gives
\begin{equation}
E_n-E'_n=\bigl(q(X_0)_{k-1}-X_0\bigr)-\bigl(q(R_n)_{k-1}-R_n\bigr).
\end{equation}
Both bracketed terms lie in $(-q/2,q/2]$, which yields
$|E_n-E'_n|<q$ for every input sequence. AMA5 implements rounding
with a one-operation delay. Swapping operands makes the carry read the fresh
top bit, so that orientation is state-blind and has the same twin.

\section{Numerical checks of the theorems}

\paragraph{Theorem 2 (law-optimal state-blind design)}
We searched all $g\in\{-2,\ldots,2\}^q$ for 132 random residue laws
with $q\le8$. In every case, a threshold quantizer reached the
minimum to $10^{-15}$. Under the unsigned and signed product-residue laws,
the best threshold
improves diffusion over round-half-up by at most 1.5\% for every $t\le8$.

\paragraph{Theorem 4 (exact moments)}
For every design, we compared diffusion from the Poisson equation with
$(\Var E_{8192}-\Var E_{4096})/4096$ from the moment recursion, using a
relative tolerance of $10^{-6}$. Direct simulations of wide accumulators
checked Lemma~2 on uniform, skewed, constant, ramp, and random-walk operand
streams. For every trace, the difference between the two sides of the
carry-conservation identity was zero. The gap in Theorem~1 remained
below $q$ for all 15 state-blind 8-bit cells. This supports the
pathwise bound under each tested input stream.

The single-class condition matters because random entry into recurrent
classes with different drifts makes variance grow quadratically.

\paragraph{Theorems 5 and 6 (bounds)}
We tested Eq.~(6) on 3000 random low-part machines with $k\le4$, at most
$2q$ states, carries in $\{-2,\ldots,3\}$, and arbitrary state encodings.
No inequality was violated. Rounding met all three bounds to
$2\times10^{-16}$ for every $k\le5$ and $s\le k$.
The 143 census design orientations also satisfied the entropy bound.
The law-dependent bound held for 1500 random machines with random
laws. It also held for all 226 design--law pairs under the product
laws, with equality for 11 pairs.

\paragraph{Relation to causal source coding}
Results for causal and finite-state source codes show that memory cannot
improve memoryless quantization of a memoryless source. Theorem~5
measures the resource in circuit registers and bounds accumulated
integer error rather than individual-sample distortion. It holds
for any logic, state encoding, or carry range. At fractional $h$, only the
first, cubic bound of Eq.~(6) can be tight (Proposition~1).

\section{Census details}

Operand orientation affects classification more often than diffusion.
Swapping the inputs changes diffusion only for \texttt{add8u\_0UK} and
\texttt{add8u\_2J3}. It changes state-blind status for 13 of 54
library cells. The state-blind counts are 26 as wired and 27 with
inputs swapped. In at least one orientation, 33 cells are
state-blind. Both orientations appear in the study records.

The closest cells to the entropy bound are \texttt{add8u\_6S4} and
\texttt{add8u\_6SM}, with the former reaching 3.91 of 4.20 bits, or 93\% of the
bound. The latter reaches 0.92 of 0.98 bits, or 94\%.
Among state-aware library cells, the recurrent-state count exceeds $2^h$
by a median factor of $2^{0.78}$.

Compensation changes how well single-addition MSE predicts accumulated MSE
at $n=256$. After optimal constant compensation, Kendall $\tau$ is 0.90
for 8-bit cells, 0.97 for 12-bit cells, and 1.00 for 16-bit cells. Without
compensation, the corresponding values are 0.54, 0.57, and 0.36. Among
434 ordered 8-bit pairs, 99 reverse. The single-addition mean error of
\texttt{add8u\_4M7} is
$-0.5$, but its accumulated drift is $-3.5$ per operation. After
compensation, its diffusion equals that of carry-cut cells
\texttt{add8u\_6QU} and \texttt{add8u\_6FT}, despite their different
single-addition MSEs.

\section{Hardware details}

\paragraph{Synthesis and mapping flow}
We run Yosys and ABC binaries from the pinned OSS CAD Suite
build on Windows; tool-build provenance and executable paths are recorded
in \texttt{asic\_flow.py}. Yosys runs \texttt{synth -noabc} and
\texttt{dffunmap}, after which standalone ABC maps the generic netlist to
Nangate45 X1 cells. Each of three structural pre-optimizations is followed
by six delay-relaxation settings of ABC's \texttt{\&nf} technology mapper,
giving the area--delay points used in the comparison. We apply no retiming
or sequential optimization, and the exact library cell
\texttt{add8u\_1HG} serves as a synthesis control and is excluded from all
counts.

The Python netlist simulator compares each register's next state
with the reference model on every cycle of a 4000-operation MAC trace
cleared every 144 operations. The combinational equivalence proof matches
registers by name, whose correspondence is fixed by construction. A
zero-delay energy estimate is about 23\% below the unit-delay estimate,
without changing the conclusions.

Because high operand bits do not affect the signed error, an exact upper
adder can replace the corresponding output logic of a 12- or 16-bit cell.
The mapper can restructure either implementation;
for 8-bit cells at $k=8$, we use the original RTL for all eight bits.

\paragraph{Quantizer-twin extraction}
The quantizer twins implement truncation, round-half-up, or truncation with
a carry-in on each operation, the last case applying to the two designs
with constant carry. We also matched each of the 1440 extracted netlists
to its twin's reference model in simulation.

Extraction did not increase delay or energy, and any reported area
increase stayed within 0.01~$\mu$m$^2$ of the mapped-cell area rounding. The
maximum area reductions were 29.4\% for unsigned and 26.0\% for signed
designs, with medians of 8.8\% and 7.8\%, respectively. Maximum energy
reductions were 25.7\% and 21.3\%, with medians of 7.0\% and 5.5\%.

\paragraph{Error measures}
Consecutive trace products are independent, so Theorem~4 applies exactly to
their residue law. Finite-length MSE and diffusion give the same comparisons
when differences below 1\% count as ties because a design and its twin differ only
by a bounded readout term, whose relative effect is $O(1/n)$. The unsigned
product residues are nonuniform: their least significant bit is one with
probability $1/4$.

\paragraph{Comparison per metric}
The per-metric result persists after excluding registers for 87 of the 89
designs. The low-part logic of \texttt{add16u\_0NL} for unsigned operands
and \texttt{add12u\_4NT} for signed operands is slightly smaller than
reference logic, but the additional registers outweigh this advantage.

\paragraph{Joint and iso-delay comparison}
Under uniform residues, exact references dominate all but 2 unsigned and 3
signed mapped points, all from state-aware designs: AMA2 at
$k=2,3$ and \texttt{add12u\_4NT}. With the product-residue law, exceptions
increase to 6 unsigned points from 3 state-aware designs. The signed
exceptions number 11 of 1602 points and come from 8 designs.

At a common delay limit, three signed designs each have one mapped point
below every qualifying reference in both area and energy.
Table~\ref{tab:isodelay} gives these exceptions. Two unsigned AMA5 points
at $k=5$ and one signed AMA2 point at $k=3$ are faster than every equally
accurate reference.

\begin{table}[t]
\centering
\caption{Area and energy differences for the signed designs that outperform
all exact references meeting the same delay limit. Negative values denote
lower cost than the corresponding reference.}
\label{tab:isodelay}
\begin{tabular}{lrr}
\toprule
Design & Area & Energy\\
\midrule
\texttt{add8u\_2LL} & $-6.4\%$ & $-7.3\%$\\
\texttt{add8u\_2J3} & $-5.1\%$ & $-4.5\%$\\
\texttt{add12u\_0C9} & $-0.9\%$ & $-0.6\%$\\
\bottomrule
\end{tabular}
\end{table}

In Table~\ref{tab:isodelay} order, the three designs have the same
diffusion under the signed product-residue law as \texttt{round\_5},
\texttt{round\_5}, and \texttt{round\_3}, respectively, to fourteen
digits. AMA5 is pathwise equivalent to rounding.
Truncation has higher diffusion under this asymmetric law and leaves the
reference set. The fastest remaining \texttt{round\_5} mapping takes
428.7~ps, compared with 411.3~ps for \texttt{add8u\_2LL}. Their
difference of about 4\% lies within implementation scatter.
The 18 mappings per design span the area--delay trade-off used for this
delay-constrained comparison.

\paragraph{Compensation cost}
Clearing to a compensation constant changes each register bit's clear logic
from forcing 0 to forcing a fixed value. Under the unsigned product-residue law, we
used $-\mathrm{round}(\E E_{144})$ for \texttt{round\_4},
\texttt{trunc\_4}, and \texttt{add8u\_6QU}. Their best mapped areas changed
by $+5.9$, $0.0$, and $-6.6~\mu$m$^2$, respectively. The corresponding
percentage changes were $+3.2\%$, $0\%$, and $-3.3\%$. These changes are
within the mapping heuristic's observed variation.

\section{Post-route and FPGA details}

\paragraph{Place-and-route flow}
We run OpenROAD-flow-scripts from the \texttt{openroad/orfs} Docker image
with its Nangate45 platform. The flow is identical for every design: ORFS
Yosys synthesis, floorplanning, I/O placement, timing-driven global
placement, detailed placement, clock-tree synthesis, timing repair, global
and detailed routing, and RC extraction. Ports are not buffered. The
default configuration uses 50\% core utilization and ORFS's speed-oriented
ABC script. The three other configurations use the area-oriented ABC
script or 40\% or 65\% core utilization. Each of the 115 designs and the 40
pruned state-blind designs is routed for both operand types at 1.5~ns and
0.55~ns under all four configurations, for 2480 routed accumulators. The
script is \texttt{pnr\_flow.py}, and the analysis is
\texttt{pnr\_analysis.py}.

\paragraph{Gate-level check and power}
A Verilog model of each Nangate45 cell is generated from its Liberty
function. Each routed netlist is simulated with Icarus Verilog on 2000
products of uniform 8-bit operands, cleared every 144 operations, at the
clock period of its run. Its outputs are compared at every cycle with the
same Python reference model used at synthesis level; for a pruned design,
the reference model is its quantizer twin. All 2480 netlists matched.
OpenSTA annotates the resulting VCD and the routed SPEF and reports total
power, which we convert to energy per accumulation by multiplying by the
clock period. Because the simulation has zero delay, glitch power is
omitted.

\paragraph{Implementation scatter}
The relative range of a design's routed area over the four
configurations has a median of 19--21\% and reaches 47\%.
Most of this range comes from the ABC script. The best-of-four
comparison removes most of this scatter. The
remaining noise is visible in designs that are pathwise within $q$ of
rounding. Over $k=2,\ldots,8$, both operand types, and both clocks, LOA and
AMA5 differ from \texttt{round\_$k$} by $-2.6\%$ to $+38\%$ in area. Their
energy differs by down to $-8.4\%$, and the median difference is $+10\%$
for area and $+9\%$ for energy.

\paragraph{Pruned designs}
For pruned state-blind RTL, the best routed area relative to the
unpruned design has a median of $-5.4\%$ (unsigned) and $-7.8\%$
(signed) at 1.5~ns. For a few unsigned designs, resynthesis produced
a larger netlist, up to $+17\%$. This is a property of the synthesis
heuristic: Corollary~1 removes cells from a given netlist, whereas the
flow resynthesizes the pruned RTL from scratch.

\paragraph{All comparisons}
Table~\ref{tab:pnrfull} lists every post-route comparison by reference
set, residue law, operand type, and clock. Not counting the pruned
netlists, or using only the default configuration, lowers the medians and
admits more designs with small negative penalties. Most are state-blind,
so Corollary~1 gives an equally cheap twin. Without pruned references,
three state-aware designs are cheaper in both area and energy, all at
1.5~ns. They are \texttt{add12u\_4NT} (signed, product-residue law, $-3.5\%$ and
$-3.9\%$), LOA with $k=2$ (unsigned, product-residue law, $-1.3\%$ and $-2.8\%$),
and \texttt{add8u\_6SM} (unsigned, both laws, $-5.3\%$ and $-4.2\%$; $-5.6\%$
and $-6.0\%$ with the default configuration only). The \texttt{add8u\_6SM}
advantage disappears once pruned state-blind netlists, which compute
the same rounding function, are counted. This indicates that ORFS
synthesizes the plain \texttt{round\_$k$} RTL suboptimally at this point.
For \texttt{add12u\_4NT}, the synthesis-level logic excluding registers was
already slightly smaller than that of the reference (Section~S5).

\begin{table}[t]
\centering
\caption{Post-route penalties (\%) relative to the cheapest equally accurate reference: medians with minima in parentheses, by reference set, diffusion law, operand type (u/s), and clock (ns).}
\label{tab:pnrfull}
\scriptsize
\setlength{\tabcolsep}{2.5pt}
\begin{tabular}{llllrrrr}
\toprule
Comparison & Law & Op. & Clock & Blind area & Blind energy & Aware area & Aware energy\\
\midrule
best of 4, with twins & uniform & u & 1.5 & 14.1 ($-1.4$) & 10.1 ($-3.5$) & 17.0 (5.0) & 12.5 ($-4.1$)\\
best of 4, with twins & uniform & u & 0.55 & 14.1 ($-0.7$) & 11.3 ($-4.9$) & 19.0 (7.4) & 15.7 (2.1)\\
best of 4, with twins & uniform & s & 1.5 & 11.3 (0.6) & 13.3 (0.5) & 12.7 (0.8) & 14.3 ($-1.2$)\\
best of 4, with twins & uniform & s & 0.55 & 11.7 (2.4) & 11.8 (2.4) & 16.6 (4.5) & 17.2 (6.3)\\
best of 4, with twins & product & u & 1.5 & 14.1 ($-1.4$) & 9.2 ($-3.5$) & 13.5 ($-1.3$) & 9.9 ($-4.1$)\\
best of 4, with twins & product & u & 0.55 & 13.5 ($-0.7$) & 10.9 ($-4.9$) & 17.1 (0.5) & 13.9 ($-1.4$)\\
best of 4, with twins & product & s & 1.5 & 10.2 (0.6) & 13.3 (0.5) & 11.7 ($-3.2$) & 13.9 ($-3.9$)\\
best of 4, with twins & product & s & 0.55 & 11.1 (2.4) & 11.8 (2.4) & 15.7 (1.6) & 15.4 ($-1.8$)\\
best of 4, no twins & uniform & u & 1.5 & 11.1 ($-4.6$) & 8.0 ($-5.4$) & 8.9 ($-5.3$) & 9.5 ($-7.6$)\\
best of 4, no twins & uniform & u & 0.55 & 12.0 ($-4.2$) & 7.9 ($-6.7$) & 12.5 (2.1) & 14.0 (1.5)\\
best of 4, no twins & uniform & s & 1.5 & 11.3 ($-1.9$) & 10.1 ($-6.8$) & 11.7 (0.1) & 10.2 ($-4.0$)\\
best of 4, no twins & uniform & s & 0.55 & 11.7 (1.6) & 10.2 ($-3.2$) & 16.1 (3.7) & 16.1 (1.3)\\
best of 4, no twins & product & u & 1.5 & 8.7 ($-4.6$) & 7.2 ($-5.4$) & 8.5 ($-5.3$) & 8.1 ($-7.6$)\\
best of 4, no twins & product & u & 0.55 & 8.2 ($-4.2$) & 7.4 ($-6.7$) & 12.0 (0.5) & 11.2 ($-1.4$)\\
best of 4, no twins & product & s & 1.5 & 9.9 ($-1.9$) & 9.7 ($-6.8$) & 10.5 ($-3.5$) & 9.5 ($-4.0$)\\
best of 4, no twins & product & s & 0.55 & 10.6 (1.6) & 9.4 ($-3.2$) & 13.5 (1.6) & 10.6 ($-1.8$)\\
default only, no twins & uniform & u & 1.5 & 11.5 ($-5.0$) & 8.3 ($-6.3$) & 8.5 ($-5.6$) & 10.6 ($-6.0$)\\
default only, no twins & uniform & u & 0.55 & 12.2 ($-4.5$) & 8.2 ($-6.7$) & 12.1 (1.9) & 12.8 (2.0)\\
default only, no twins & uniform & s & 1.5 & 11.2 ($-1.9$) & 9.7 ($-6.6$) & 11.4 (0.1) & 10.5 ($-4.2$)\\
default only, no twins & uniform & s & 0.55 & 11.8 (1.6) & 11.0 ($-3.2$) & 13.2 (3.9) & 12.7 (1.3)\\
default only, no twins & product & u & 1.5 & 8.7 ($-5.0$) & 7.1 ($-6.3$) & 8.5 ($-5.6$) & 7.8 ($-6.0$)\\
default only, no twins & product & u & 0.55 & 8.3 ($-4.5$) & 7.5 ($-6.7$) & 10.9 (0.5) & 10.0 ($-1.3$)\\
default only, no twins & product & s & 1.5 & 9.9 ($-1.9$) & 9.2 ($-6.6$) & 10.2 ($-3.3$) & 9.0 ($-4.2$)\\
default only, no twins & product & s & 0.55 & 10.7 (1.6) & 10.5 ($-3.2$) & 12.1 (1.4) & 10.4 ($-1.8$)\\
\bottomrule
\end{tabular}
\end{table}

\begin{table}[t]
\centering
\caption{Routed registered $8\times8$ signed MAC (default configuration, 2~ns clock): standard-cell area and change relative to the exact MAC (\%).}
\label{tab:mac}
\footnotesize
\begin{tabular}{lrr}
\toprule
Accumulator & Area ($\mu$m$^2$) & vs.\ exact\\
\midrule
\texttt{exact} & 924.1 & 0.0\\
\texttt{round\_5} & 823.8 & $-10.9$\\
\texttt{round\_6} & 842.4 & $-8.8$\\
\texttt{round\_7} & 773.0 & $-16.4$\\
\texttt{round\_8} & 753.6 & $-18.5$\\
\texttt{trunc\_4} & 869.6 & $-5.9$\\
\texttt{trunc\_6} & 798.3 & $-13.6$\\
\texttt{trunc\_7} & 774.3 & $-16.2$\\
\texttt{add8u\_0E2} & 809.7 & $-12.4$\\
\texttt{add8u\_6QU} & 842.2 & $-8.9$\\
\texttt{add8u\_2LL} & 825.4 & $-10.7$\\
\texttt{loa\_5} & 878.9 & $-4.9$\\
\bottomrule
\end{tabular}
\end{table}


\paragraph{FPGA}
Vivado 2022.2 implements each signed design out of context on
\texttt{xck26-sfvc784-2LV-c} with a 2~ns clock. The flow runs
\texttt{synth\_design -flatten\_hierarchy full}, \texttt{opt\_design},
\texttt{place\_design}, and \texttt{route\_design}. All 155 designs,
including the pruned ones, met the clock. The exact accumulator uses 24
LUTs, 24 registers, and three \texttt{CARRY8} blocks. \texttt{round\_8}
uses 16 LUTs and 16 registers. The script is \texttt{fpga\_flow.py}, and the
analysis is \texttt{fpga\_analysis.py}.

\paragraph{MAC}
Table~\ref{tab:mac} gives the routed area of a registered signed
$8\times8$ MAC built around each accumulator in the CNN equivalence set.
It uses the default configuration only.

\section{CNN details}

\paragraph{Setup}
The second training seed reached 91.49\%, and activation scales were
calibrated with 500 training images. Across the test set, the largest exact
accumulator magnitude was 180{,}578, well within the 24-bit range. We
measured the product-residue law for compensation on 40 calibration images,
sampling 8 from each batch of 100.

At top-1 accuracy near 92\%, the 95\% confidence half-width on 1000 images
is 1.7 points, so we also compare designs image by image.

\paragraph{Model accuracy}
We check Theorem~4 in sample, using each layer's residue law measured on the
test images. Across 494 layer--design pairs, the median ratio of predicted
to measured uncompensated MSE is 0.978. The interquartile range is
0.961--0.989, the 5--95\% range is 0.935--1.002, and the maximum is 2.0.
Under the uniform law, the median ratio is 3.4, the 5--95\% range is
0.52--44, and the extremes are 0.29 and 97.

\paragraph{Second seed}
Table~\ref{tab:secondseed} lists calibrated accuracy for eleven
representative designs in the second network. Their values span
90.7--91.3\%, compared with 91.0\% for the exact accumulator.
Predictions from \texttt{add8u\_6QU} and \texttt{trunc\_4} agree on 99.4\%
of images.
\begin{table}[t]
\centering
\caption{Calibrated classification accuracy for the second training seed.
The paired columns list designs from the same evaluation, but they do not
represent paired comparisons between rows.}
\label{tab:secondseed}
\begin{tabular}{lr@{\qquad}lr}
\toprule
Design & Cal. top-1 (\%) & Design & Cal. top-1 (\%)\\
\midrule
exact & 91.0 & \texttt{round\_4} & 91.0\\
\texttt{round\_6} & 91.0 & \texttt{trunc\_4} & 90.9\\
\texttt{trunc\_6} & 90.7 & \texttt{add8u\_6QU} & 91.1\\
ETA-I, $k=5$ & 91.3 & LOA, $k=5$ & 91.1\\
\texttt{add8u\_6SM} & 91.0 & \texttt{add8u\_2LL} & 91.3\\
AMA2, $k=3$ & 91.2 & \texttt{add8u\_0PL} & 91.0\\
\bottomrule
\end{tabular}
\end{table}
\paragraph{Requantization headroom}
Per output channel, requantization to INT8 after each layer discards at
least 7.4 accumulator least significant bits, with per-layer medians of
8.5--10.6 bits.

\paragraph{Equivalence testing}
The one-point margin is an application-level tolerance selected after the
initial 1000-image analysis. We then ran a paired two one-sided test (TOST)
at level 0.05. It used per-image correctness differences from the exact
accumulator and a $\pm1$-point margin. The 10{,}000-image evaluation was
also added after the first analysis of 1000 images.

Under the same 90\% paired intervals but a half-point margin,
\texttt{round\_5}, \texttt{round\_6}, \texttt{round\_7},
\texttt{trunc\_4}, \texttt{add8u\_0E2}, \texttt{add8u\_6QU},
\texttt{add8u\_2LL}, and LOA at $k=5$ remain equivalent.
\texttt{round\_8} and \texttt{trunc\_6} do not. The former has an interval of
$[-0.784,-0.296]$ percentage points, so its classification as
equivalent at 33\% less area depends on the one-point tolerance
(\path{cnn_cal_10k_tost_half_point.json}).

On 1000 images, most evaluated designs are equivalent to exact.
The exceptions are \texttt{trunc\_7} and \texttt{trunc\_8}, whose
differences are $-2.1$ and $-54.4$ points. The result for
\texttt{round\_8} is inconclusive. The
last difference is $-0.6$ points, with a 90\% interval of $[-1.4,+0.2]$.
At this sample size, \texttt{round\_7} is the cheapest equivalent
accumulator. Its difference is $+0.1$ points, with interval
$[-0.5,+0.7]$. Its area is 183.8~$\mu$m$^2$, or 28\% less than exact.

On 10{,}000 images, \texttt{round\_7} loses 0.11 points, with interval
$[-0.28,+0.06]$. The \texttt{round\_8} loss is 0.54 points, with
interval $[-0.78,-0.30]$. The cells
\texttt{add8u\_0E2}, \texttt{add8u\_6QU}, \texttt{add8u\_2LL}, and LOA at
$k=5$ remain within 0.12 points. Quantizer-twin predictions agree on 99.6\%
of images for \texttt{add8u\_6QU} and \texttt{trunc\_4}, 99.6\% for LOA and
\texttt{round\_5}, and 98.9\% for \texttt{add8u\_0E2} and
\texttt{round\_7} among the 10{,}000 images. The remaining
disagreements reflect the bounded readout term and small differences
between calibrated preloads.

\paragraph{Energy on CNN operands}
We re-simulated every mapped signed netlist on 28 serial accumulations of
real products. The products came from four test images and two layers, with
$L=144$ and $576$, for 10{,}080 operations in total.

\section{CIFAR-100 replication details}

\paragraph{Protocol}
The protocol was written before the CIFAR-100 network was trained and is
included in the public artifact under \texttt{paper/integration\_carry/qa/}. Its SHA-256 hash was
recorded in the revision log at that time. It fixes the network, training
recipe, quantization, compensation, test set, the twelve designs, the two
area rankings, and hypotheses H1--H4. The network was trained for 50
epochs with seed~0 and reached 67.49\% float test accuracy. The analysis
script is \texttt{cifar100\_analysis.py}, and its output is
\texttt{cifar100\_analysis.json}.

\paragraph{Equivalence}
At the secondary half-point margin, the same nine designs are equivalent.
McNemar tests against exact give $p<10^{-6}$ for \texttt{round\_8},
\texttt{trunc\_6}, and \texttt{trunc\_7}, and $p\ge0.11$ for every other
design.

\paragraph{Prediction agreement}
On CIFAR-100, prediction agreement with exact varies across
references. It is 97.4\% for \texttt{round\_5} and 96.2\% for
\texttt{round\_6}. It is 93.6\% for \texttt{round\_7} and 87.1\%
for \texttt{round\_8}. For \texttt{trunc\_4}, it is 96.4\%.
The corresponding CIFAR-10 values on 10{,}000 images
are 97.4--99.4\%. \texttt{add8u\_0E2} agrees with \texttt{round\_7} on
95.0\% of images and with every other reference on at most 93.7\%;
\texttt{add8u\_6QU} agrees with \texttt{trunc\_4} on 98.1\% and with
every other reference on at most 96.4\%; LOA agrees with \texttt{round\_5}
on 97.9\% and with every other reference on at most 97.5\%.

\paragraph{Variance prediction}
For H2, each layer's prediction uses Theorem~4 under the pooled
residue law. It is $\mathrm{Var}+(\E E_K+p)^2$, where $p$ is the applied
preload. Over 209 layer--design pairs, the median ratio is
0.876. Taking the variance alone, the per-design median ratio is 0.94--0.99
for \texttt{round\_5}--\texttt{round\_8}, LOA, \texttt{add8u\_0E2}, and
\texttt{add8u\_2LL}, but 0.20--0.25 for \texttt{trunc\_4},
\texttt{trunc\_6}, \texttt{trunc\_7}, and \texttt{add8u\_6QU}. The
CIFAR-10 10{,}000-image run gives the same split (0.94--0.99 and
0.22--0.27).

The post hoc diagnostic \texttt{mse\_sparsity\_diag.py} computes twin
error exactly from products of 24 test images in each convolution
layer. It covers \texttt{trunc\_4}, \texttt{trunc\_6}, and
\texttt{round\_5}.
Using per-output-channel residue laws changes the ratios little (median
0.26 for \texttt{trunc\_4}). If $N$ is the number of nonzero products in
an output, and the nonzero residues are independent with per-operation
error mean $\mu$ and variance $\sigma^2$, then
$\mathrm{Var}(E)=\E N\,\sigma^2+\mathrm{Var}(N)\,\mu^2$. With this
conditioning, the median ratios are 1.00 for all three designs on both
datasets, and every layer lies in 0.95--1.04. The second term is negligible
for rounding, where $\mu\approx0$, which explains why rounding-type
designs were already predicted well.

\section{FIR details}

For each $L$, we estimate the product-residue law and compensation
constants from 20{,}000 windows in images 1--12. We measure SNR on another
20{,}000 held-out windows from images 13--24. We form 95\%
bootstrap intervals over windows with 1000 resamples, using the same
17 references and 18 approximate designs as the CNN evaluation.

The bounded readout term has a larger effect on shorter filters. At
$L=256$, \texttt{add8u\_0TP} reaches exactly the 46.64~dB of
\texttt{trunc\_5}. At $L=64$, the largest difference from a quantizer twin
is 0.5~dB, for LOA at $k=5$.

Theorem~4 predicts error MSE with median predicted-to-measured ratios
of 0.98 and 0.96. Their respective interquartile ranges are
0.88--1.14 and 0.78--1.09. Neighboring pixels make successive
products dependent, so individual errors can be larger. The model
underestimates AMA2 and \texttt{add8u\_0PL} by up to 17 and 1.9
times. It misestimates single-bit rounding and eight-bit truncation
by 0.56 and 3.3.

For the hardware comparison, we re-simulate every signed netlist on
40 serial FIR accumulations of held-out products per length.
Under the strict rule that a reference must have at least the design's
measured SNR, area penalties range from $+2.4\%$ to $+22.0\%$.
Energy penalties range from $+4.0\%$ to $+28.3\%$.
If overlapping SNR intervals count as equal
accuracy, the minimum area penalty at $L=256$ becomes $+4.6\%$. The energy
penalty range then becomes $+8.4\%$ to $+28.3\%$.

\end{document}
